\documentclass[tikz,border=8pt]{standalone}
\usepackage{ctex}
\usepackage{tikz}
\usetikzlibrary{arrows.meta,positioning,fit,backgrounds}
\begin{document}
\begin{tikzpicture}[
  font=\small,
  >=Latex,
  blk/.style={draw=#1!75!black, fill=#1!10, rounded corners=2.5mm, very thick, minimum width=34mm, minimum height=11mm, align=center},
  sig/.style={->, very thick, draw=black!75},
  tag/.style={draw=black!35, rounded corners=1.5mm, fill=yellow!18, inner sep=2pt}
]
\node[blk=blue] (c1) at (0,1.8) {创新点 1\\CAN ClockIDS\\时钟倾斜建模};
\node[blk=teal] (c2) at (0,-1.8) {创新点 2\\以太网环形存储\\时间窗关联};
\node[blk=purple] (c3) at (5.8,0) {创新点 3\\CAN-CNN 64x9\\载荷+时序联合};
\node[blk=orange] (c4) at (11.6,0) {创新点 4\\中央处理器融合\\跨域告警解释};

\draw[sig] (c1.east) -- (c3.west);
\draw[sig] (c2.east) -- (c3.west);
\draw[sig] (c3.east) -- (c4.west);

\node[tag] at (2.8,1.0) {\texttt{real-time}};
\node[tag] at (2.7,-1.0) {\texttt{O(1) write}};
\node[tag] at (8.7,1.0) {\texttt{64x9 compact}};
\node[tag] at (8.7,-1.0) {\texttt{interpretable}};

\begin{scope}[on background layer]
  \node[draw=black!25, rounded corners=3mm, line width=0.8pt, fill=white,
  fit=(c1)(c2)(c3)(c4), inner sep=8mm, label={[yshift=2mm]\large\bfseries 核心创新点总览}] {};
\end{scope}
\end{tikzpicture}
\end{document}
